\chapter{Estado del arte}\label{Cap4}
Hay varios métodos existentes que están disponibles en el grado industrial para limpiar un panel solar que contiene las problemáticas de terreno geográfico, ambientales y su área de aplicación. Dependiendo de estos factores existe una base de comparación en cuestión de uso, desempeño, viabilidad, etc. Una de estas soluciones es un robot automatizado que pueda reducir costos y tiempos empleados para limpiar paneles solares [6].\\

Como opción más innovadora se tiene el Solarbrush UAV Robot: el robot UAV solarbrush, hasta ahora, desempeña el papel de limpiar los paneles solares de una manera única, dado a que el robot está compuesto de 4 hélices y una cola que tiene la tarea específica de limpiar el panel siendo el único elemento que hace contacto con la superficie del panel.\\

La única parte del dron que hace contacto con la superficie del panel fotovoltaico es la cola, la cual, según sus desarrolladores, es la gran ventaja de usar este robot, dado a que esta propiedad ayuda que la probabilidad de dañar el panel fotovoltaico decrece.
\begin{savenotes}
\begin{figure}[H]
	\centering
	\includegraphics[width=0.7\textwidth,keepaspectratio]{uav}
	\caption{Robot solarbrush UAV.}
	\label{figuav}
\end{figure}
\end{savenotes}
Ecoppia E4: Ecopia E4, como se ven en la Figura \ref{figecoppia}, es un robot totalmente autónomo que usa como agente limpiador una microfibra que cepilla de manera limpia el polvo de los paneles. Está diseñada para limpiar grandes cantidades de paneles localizados en áreas secas y de terracería. El robot usa un módulo fotovoltaico que también cuenta con almacenamiento de energía para poder realizar la limpieza durante la noche, y el posicionamiento de la celda fotovoltaica en la parte inferior del riel le permite que sea limpiada por el robot mismo y el robot tiene implementada inteligencia para desempeñar su trabajo de limpieza acorde a las condiciones climáticas.\\

EL movimiento del robot es ejecutado a través de un riel guía que está implementado sobre la esquina superior e inferior del panel fotovoltaico, además de que este robot tiene la adición de ser controlado de forma remota, donde los comandos de limpieza pueden ser obtenidos por cualquier equipo conectado a internet.
\begin{savenotes}
\begin{figure}[H]
	\centering
	\includegraphics[width=0.7\textwidth,keepaspectratio]{ecoppia}
	\caption{Robot Ecoppia E4.}
	\label{figecoppia}
\end{figure}
\end{savenotes}
Método de rociadores para limpiar paneles fotovoltaicos, Estos rociadores son usados frecuentemente en áreas secas para mantener los paneles limpios sin obstrucciones de polvo, implementando un método de limpieza rociando cierto tipo de agua con jabón que limpia el panel por un costo relativamente bajo.\\

Este método consiste en la filtración del agua y un sistema dispensador de jabón, aunque éste es un relativamente buen sistema de limpieza, hay algunas desventajas recurrentes para considerar que no son tan convenientes, la primera es el gasto de agua, dependiendo de cuantas veces se limpien los paneles por día, otra es la complicación de limpiar los paneles de manera correcta debido a un sobre acumulamiento de tierra o polvo, haciendo que sea más difícil para los rociadores el remover la suciedad debido a que no tiene un agente que talla los paneles. Por lo tanto, para resolver estos problemas se le da un monitoreo regular al sistema haciendo que cueste más tiempo invertido en revisar su correcta limpieza y en que puede llegar a costar mucho dinero por reaprovisionamiento de agua y jabón.
\begin{savenotes}
\begin{figure}[H]
	\centering
	\includegraphics[width=0.7\textwidth,keepaspectratio]{ecoppia}
	\caption{Rociadores usados para limpiar paneles fotovoltaicos.}
	\label{figrociador}
\end{figure}
\end{savenotes}
Así como el mencionado Ecoppia E4, el NOMADD (por sus siglas en inglés: NO Water Mechanical Automated Dusting Device), mostrado en la Figura \ref{fignomadd}, es un robot diseñado específicamente para remover la acumulación de polvo de arreglos de paneles fotovoltaicos del desierto sin el uso de agua.\\

Contiene la propiedad de limpiar paneles sin agua, además de tener su propia celda solar con la adición de ser totalmente operable remotamente y automatizado, con la capacidad de seguir trabajando aún con condiciones adversas debido a su robusto diseño.
\begin{savenotes}
\begin{figure}[H]
	\centering
	\includegraphics[width=0.7\textwidth,keepaspectratio]{nomadd}
	\caption{Robot NOMADD Limpiador.}
	\label{fignomadd}
\end{figure}
\end{savenotes}
